\section[向量丛的族]{向量丛的族（11 月 27 日）}

\srcpage{42}

\subsection{例子：Hodge--Tate 周期映射}

我们已经看到，对 $C$ 上的 Fargues--Fontaine 曲线，等晶体给出向量丛，而所有向量丛都由此得到。我们想先看一个例子，说明这两个概念在族中如何不同。

考虑椭圆曲线的模空间 $\mathcal{M}_{\mathrm{ell}}/\mathbb{Z}$。它是 Deligne--Mumford 叠；不过若施加一点辅助水平结构，它会变成仿射概形。这个区别对我们的讨论无关紧要。

在特征 $p$ 中，我们有两个分层，分别参数化普通（ordinary）与超奇异（supersingular）椭圆曲线：
\[
\mathcal{M}_{\mathrm{ell},\Fp}
=
\underbrace{\mathcal{M}^{\mathrm{ord}}_{\mathrm{ell},\Fp}}_{\text{开}}
\;\cup\;
\underbrace{\mathcal{M}^{ss}_{\mathrm{ell},\Fp}}_{\text{闭（有限）}}.
\]
说明 $E/k=\overline{\Fp}$ 为普通的一种方式，是看 $H^1_{\mathrm{crys}}(E/W(k))[1/p]\in\Isoc_{\Qp}$。它是一个等晶体，或为
\[
\underbrace{\Bigl(\breve{\mathbb{Q}}_p^2,
\begin{pmatrix}p&\\&1\end{pmatrix}\sigma\Bigr)}_{\text{普通（斜率 0、1）}}
\qquad\text{或}\qquad
\underbrace{\Bigl(\breve{\mathbb{Q}}_p^2,
\begin{pmatrix}1&\\&p\end{pmatrix}\sigma\Bigr)}_{\text{超奇异（斜率 $1/2$）}}.
\]
因此，在 $\mathcal{M}_{\mathrm{ell},\Fp}$ 上我们有一族从普通退化到超奇异的等晶体。（为了理解「等晶体的族」，考虑 $\mathcal{M}_{\mathrm{ell},\Fp}$ 的完满化。在一个（完满的）仿射开集上，取万有族的晶体上同调。）然而，在研究 FF 曲线上的向量丛时，这幅图景反过来了。

考虑 $\mathcal{M}_{\mathrm{ell},\Fp}\subset\mathcal{M}^{\wedge}_{\mathrm{ell},\mathbb{Z}}\big|_{\mathbb{Z}_p^{\wedge}}$，其中后者被视为 $\mathbb{Z}_p^{\wedge}=\cO_{C_p}$ 上的形式概形。它有 adic 一般纤维 $\mathcal{M}^{\mathrm{ad}}_{\mathrm{ell},C_p}$。存在特殊化映射
\[
\mathrm{sp}\colon\bigl|\mathcal{M}^{\mathrm{ad}}_{\mathrm{ell},C_p}\bigr|\to\bigl|\mathcal{M}_{\mathrm{ell},\overline{\Fp}}\bigr|,
\]
它连续，并通过拉回诱导出 $\mathcal{M}_{\mathrm{ell},C_p}$ 上类似的普通—超奇异分层。

另一方面，$\mathcal{M}_{\mathrm{ell}}$ 有一个带无穷级水平结构的版本
\[
\mathcal{M}_{\mathrm{ell},C_p,\infty}\sim\varprojlim_{m}\mathcal{M}_{\mathrm{ell},C_p,p^m},
\]
其中 $\mathcal{M}_{\mathrm{ell},C_p,p^m}$ 参数化椭圆曲线 $E$ 连同水平结构 $E[p^m]\cong(\mathbb{Z}/p^m)^2$。因此在左边我们有完全平凡化 $E[p^\infty]\cong(\mathbb{Q}_p/\mathbb{Z}_p)^2$。在 [S15] 中我们证明了 $\mathcal{M}_{\mathrm{ell},C_p,p^\infty}$ 作为完满空间存在。

一个只在无穷级才出现的新现象是 Hodge--Tate 周期映射
\[
\pi_{\mathrm{HT}}\colon\mathcal{M}_{\mathrm{ell},C_p,p^\infty}\to\mathbb{P}^1_{C_p}.
\]
简单回顾一下：$E/C_p$ 有 Hodge--Tate 滤过
\[
0\to(\Lie E)(1)\to T_p(E)\otimes_{\mathbb{Z}_p}C_p\to(\Lie E^*)^*\to0.
\]
于是完全水平结构给出同构 $T_p(E)\otimes_{\mathbb{Z}_p}C_p\cong C_p^2$。

\begin{proposition}
$E$ 有普通约化，当且仅当 Hodge--Tate 滤过是 $\mathbb{Q}_p$-有理的。
\end{proposition}

要点是：若 $E$ 有普通约化，则 $T_p(E)$ 上有自然的滤过（由斜率滤过给出，即模 $p$ 的连通—étale 序列的提升）。因此 Hodge--Tate 滤过下降到 $\mathbb{Q}_p$ 上。

\srcpage{43}
\begin{center}
\figP{043-1}
\end{center}
\begin{center}
\figP{043-2}
\end{center}

有点悖谬的是，Hodge--Tate 周期映射把 $\mathcal{M}^{\mathrm{ord}}_{\mathrm{ell},C_p,p^\infty}$ 送到 $\mathbb{P}^1(\mathbb{Q}_p)$，而主要是在超奇异轨迹上才有事情发生。

本质上，$\mathcal{M}^{\mathrm{ord}}_{\mathrm{ell},C_p,p^\infty}=\pi_{\mathrm{HT}}^{-1}(\mathbb{P}^1(\mathbb{Q}_p))$ 是闭子集 $\mathbb{P}^1(\mathbb{Q}_p)\subset\mathbb{P}^1_{C_p}$ 的原像。「本质上」这个限定是必要的，因为这实际上只在秩 1 点上成立。（否则 $\mathcal{M}^{\mathrm{ord}}_{\mathrm{ell},C_p,p^\infty}$ 就会是开闭的，但 $\mathcal{M}_{\mathrm{ell},C_p,p^\infty}$ 连通。）存在某些秩 2 点，它们「刚刚好」特殊化到超奇异点，其像也落在 $\mathbb{P}^1(\mathbb{Q}_p)$ 中。

我们来放大看超奇异圆盘的原像上发生了什么。

当你走向超奇异约化圆盘的边界时，Hodge--Tate 周期映射下的像走向 $\Omega^2$ 的边界，也就是 $\mathbb{P}^1(\mathbb{Q}_p)$。在 adic 空间中，超奇异约化圆盘的边界上确实有一个秩 2 点，它映射到 $\mathbb{P}^1(\mathbb{Q}_p)$ 中。

特别地，这感觉不像复流形之间的映射。在 $\Omega^2$ 上纤维是投射有限集，而在 $\mathbb{P}^1(\mathbb{Q}_p)$ 上纤维是 1 维的。

\begin{remark}
这幅图景说明了无穷级的 Lubin--Tate 塔与无穷级的 Drinfeld 塔之间的同构。
\end{remark}

\srcpage{44}
\begin{center}
\figP{044}
\end{center}

\subsection{FF 曲线上向量丛的族}

用 Fargues--Fontaine 曲线上的向量丛来说，上例中的 $\mathbb{P}^1_{C_p}$ 参数化「在某个 untilt 处对平凡秩 2 向量丛的修正」。更确切地说，我们有一个 1 次除子
\[
\Spa C_p\hookrightarrow X_{C_p^\flat,\mathbb{Q}_p}.
\]
于是选取变化的直线 $C_p^2\twoheadrightarrow L$，就得到 $X_{C_p^\flat,\mathbb{Q}_p}$ 上的丛 $\mathcal{E}(L)$，由短正合列
\[
\begin{tikzcd}[column sep=small,row sep=small]
0 \arrow[r] & \mathcal{E}(L) \arrow[r] & \cO^2_{X_{C_p^\flat}} \arrow[r] \arrow[d] & L \arrow[r] & 0\\
& & C_p^2 \arrow[ur] & &
\end{tikzcd}
\]
决定。
\begin{itemize}
  \item 在 $\Omega^2\subset\mathbb{P}^1_{C_p}$ 上，$\mathcal{E}(L)=\cO(-1/2)$（它是半稳定的）。
  \item 在 $\mathbb{P}^1(\mathbb{Q}_p)\subset\mathbb{P}^1_{C_p}$ 上，$\mathcal{E}(L)\cong\cO\oplus\cO(-1)$。
\end{itemize}

现在像之前一样，固定 $E$ 为剩余域 $\Fq$ 的非阿基米德局部域，$\pi\in\cO_E$ 为一致化子。

设 $S\in\Perf_{\Fq}$ 是完满空间。则有 Fargues--Fontaine 曲线 $X_S=X_{S,E}$。设 $\mathcal{E}$ 是 $X_S$ 上的向量丛。对每个「几何点」\footnote{在与代数几何的类比中，这更像几何点处的严格 Hensel 化。adic 空间的映射是泛化的，所以不能有像为「单点」的映射。}$\bar{s}=\Spa(C,C^+)\to S$，我们可以考虑 $\mathcal{E}_{\bar{s}}/X_{\bar{s}}$。

注意：$\mathrm{VB}(X_{\bar{s}})\cong\mathrm{VB}(X_{\Spa(C,\cO_C)})$，所以可以忘掉 $C^+$（加号子环在向量丛理论中不起作用，尽管它在 étale 上同调理论中会很重要）。

我们有向量丛的分类 $\mathcal{E}_{\bar{s}}\cong\bigoplus_{\lambda\in\mathbb{Q}}\cO_{X_{\bar{s}}}(\lambda)^{n_\lambda(\bar{s})}$。于是可以构造 Newton 多边形：横轴为秩，纵轴为次数，各段的水平长度是 $\lambda\cdot n_\lambda(\bar{s})$。

\srcpage{45}
\begin{center}
\figP{045}
\end{center}

\begin{example}
例如，$\cO(-1/2)$ 的 Newton 多边形是从 $(0,0)$ 到 $(2,-1)$ 的直线段；$\cO\oplus\cO(-1)$ 的 Newton 多边形是折线 $(0,0)\to(1,-1)\to(2,-1)$。

Newton 多边形在向量丛的族中如何变化？我们定义如下序关系：
\[
P\ge P'\quad\text{当且仅当 }P\text{ 与 }P'\text{ 端点相同，且 }P\text{ 位于 }P'\text{ 之上或与之重合}.
\]
\end{example}

\begin{theorem}[Kedlaya--Liu '15]
(1) 函数 $\bar{s}\mapsto\mathrm{NP}(\mathcal{E}_{\bar{s}})$ 定义映射 $|S|\to\{\text{Newton 多边形}\}$，它是半连续的。

(2) 若 Newton 多边形为常值，则存在整体的 Harder--Narasimhan 滤过
\[
\mathcal{E}^{\ge\lambda}\subseteq\mathcal{E},
\]
由（饱和的）向量子丛给出；且每个 $\mathcal{E}^\lambda:=\mathcal{E}^{\ge\lambda}/\bigcup_{\lambda'>\lambda}\mathcal{E}^{\ge\lambda'}$ 处处是斜率 $\lambda$ 的半稳定丛。在 $S$ 上 pro-étale 局部，有同构
\[
\mathcal{E}\cong\bigoplus_{\lambda\in\mathbb{Q}}\cO_{X_S}(\lambda)^{n_\lambda}.
\]
\end{theorem}

我们将给出一个不同的（短得多的）证明，它同样依赖于菱形与 v-下降的几何。关键在于射影化的 Banach--Colmez 空间是紧合的。为解释这一点，我们需要过一遍一些技术基础。

\srcpage{46}

\subsection{分离与紧合映射}

回到 $\Perf:=\{\text{特征 }p\text{ 的完满空间}\}$ 上 v-层的设定。

我们有以下赋值判别法。

\begin{definition}
设 $f\colon\mathcal{F}\to\mathcal{G}$ 是小 v-层之间的映射。我们说：

(1) $f$ 是闭浸入，若对一切严格完全间断的 $X$ 及一切映射 $X\to\mathcal{G}$，纤维积 $\mathcal{F}\times_{\mathcal{G}}X$ 可由完满空间 $X'$ 表示，且 $X'\to X$ 是（Zariski）闭浸入。（等价地，存在泛化闭子集 $Z\subset|\mathcal{G}|$，使得 $\mathcal{F}\subset\mathcal{G}$ 是由所有满足「$|X|\to|\mathcal{G}|$ 经过 $Z$ 分解」的映射 $X\to\mathcal{G}$ 构成的子函子。）

(2) $f$ 是分离的，若 $\Delta_f$ 是闭浸入。

(3) $f$ 是紧合的，若 $f$ 分离、拟紧且泛闭。
\end{definition}

\begin{proposition}
小 v-层之间的映射 $f\colon\mathcal{F}\to\mathcal{G}$ 是分离的（相应地，紧合的），当且仅当它拟分离（相应地，qcqs），且对一切图表
\[
\begin{tikzcd}[column sep=large]
\Spa(R,R^\circ) \arrow[r] \arrow[d] & \mathcal{F} \arrow[d]\\
\Spa(R,R^+) \arrow[r] \arrow[ur,dashed] & \mathcal{G}
\end{tikzcd}
\]
以及一切 affinoid 完满空间 $\Spa(R,R^+)$，至多存在一个（相应地，恰好存在一个）虚线箭头。事实上，只需对 $(R,R^+)=(C,C^+)$ 验证，其中 $C$ 是完全代数闭域，$C^+$ 是赋值子环。
\end{proposition}

在 adic 几何中有许多合理但非拟紧的空间，例如开单位圆盘。因此我们希望有一个「不带拟紧性」的紧合概念。

\begin{definition}
小 v-层之间的映射 $f\colon\mathcal{F}\to\mathcal{G}$ 是部分紧合的，若它分离，且存在唯一的填充
\[
\begin{tikzcd}[column sep=large]
\Spa(R,R^\circ) \arrow[r] \arrow[d] & \mathcal{F} \arrow[d]\\
\Spa(R,R^+) \arrow[r] \arrow[ur,dashed,"\exists!"] & \mathcal{G}
\end{tikzcd}
\]
\end{definition}

让我们回来处理上次隐含使用的一个性质。

\begin{proposition}
设 $f\colon\mathcal{F}\to\mathcal{G}$ 拟紧。则 $f$ 作为 v-层的映射是满射，当且仅当 $|f|\colon|\mathcal{F}|\to|\mathcal{G}|$ 是满射。
\end{proposition}

\begin{proof}[证明梗概]
化归到可表情形 $\mathcal{F}=X$、$\mathcal{G}=Y$。但此时 $X\to Y$ 是 v-覆盖。（「只要一切是拟紧的，v-覆盖就是满射。」）
\end{proof}

\subsection{射影化的 Banach--Colmez 空间}

\begin{proposition}
设 $S\in\Perf_{\Fq}$，$\mathcal{E}\in\mathrm{VB}(X_S)$。则 $BC(\mathcal{E})\colon T_{/S}\mapsto H^0(X_T,\mathcal{E}|_{X_T})$ 是局部空间菱形，且在 $S$ 上部分紧合。射影化的 Banach--Colmez 空间 $(BC(\mathcal{E})\setminus\{0\})/\underline{E}^\times$ 是局部空间菱形，且在 $S$ 上紧合。
\end{proposition}

\srcpage{47}

\begin{proof}[证明梗概]
由 $\cO(1)$ 的丰沛性，当 $n,N\gg0$ 时存在满射 $\cO_{X_S}(-n)^N\twoheadrightarrow\mathcal{E}^\vee$。对偶地得到嵌入 $\mathcal{E}\hookrightarrow\cO_{X_S}(n)^N$。可以说明这给出闭浸入 $BC(\mathcal{E})\hookrightarrow BC(\cO_{X_S}(n)^N)$。局部空间菱形中的闭嵌入仍是局部空间的，因此化归到 $\mathcal{E}=\cO_{X_S}(n)^N$。部分紧合定义中赋值判别法那一部分很清楚，因为向量丛理论不依赖于 $+$ 子环。

假设 $S$ 是 qcqs 的。困难在于证明 $(BC(\mathcal{E})\setminus\{0\})/\underline{E}^\times$ 的全空间是 qcqs 的。（对概形，qcqs 部分通常容易，赋值判别法才难。在 $p$ 进几何中往往相反。）我们可以归纳地分析 $BC(\cO(n))$：
\[
0\to\cO_{X_S}(n-1)\to\cO_{X_S}(n)\to\cO_{S^\#}\to0.
\]
这给出
\[
0\to BC(\cO_{X_S}(n-1)^N)\to BC(\cO_{X_S}(n)^N)\to(\mathbb{A}^N_{S^\#})^\diamond\to0.
\]
现在目标是一个拟分离的局部空间菱形。纤维由归纳也是局部空间菱形。这个正合列在目标上 pro-étale 局部分裂，所以中间项也是局部空间的。这个论证还给出拟分离性。

困难的部分是 $(BC(\mathcal{E})\setminus\{0\})/\underline{E}^\times$ 拟紧。由于 $\cO_E^\times$ 紧，主要点是 $(BC(\mathcal{E})\setminus\{0\})/\pi^{\mathbb{Z}}$ 拟紧。这是一个纯拓扑的陈述，它由关于局部谱空间的自同构「收缩」的一般引理推出，该引理说：

设 $\gamma$ 作用在 $T$ 上，$T_0=T^\gamma$。假设：
\begin{itemize}
  \item 「$T$ 看起来像解析 adic 空间」（一个点的泛化构成全序链），且
  \item 「当 $n\to\infty$ 时，$\gamma^n$ 的作用向 $T_0$ 收缩」，
  \item 「当 $n\to-\infty$ 时，$\gamma^n$ 在 $T\setminus T_0$ 上的作用发散」。
\end{itemize}
结论：$\gamma$ 自由且不连续地作用在 $T\setminus T_0$ 上，并且 $(T\setminus T_0)/\gamma^{\mathbb{Z}}$ 是谱的，即 qcqs。证明是一些难看的点集拓扑。
\end{proof}

\subsection{定理 9.4 的证明}

设 $S\in\Perf_{\Fq}$，$\mathcal{E}\in\mathrm{VB}(X_S)$。

(1) $\mathrm{NP}(\mathcal{E})\colon\bar{s}\mapsto\mathrm{NP}(\mathcal{E}_{\bar{s}})$ 是半连续的。

初等观察：$\mathrm{NP}(\mathcal{E}_{\bar{s}})$ 是所有点 $(i,d_i)$（$i=0,1,\dots,\mathrm{rank}\,\mathcal{E}_{\bar{s}}$）的凸包，其中存在 $(\wedge^i\mathcal{E}_{\bar{s}})(-d_i)$ 的非零截面。这是用分类定理做的一个练习。

所以只要能检测出有非零截面的丛，就能检测出 NP。因此只需证明：$\mathcal{E}$ 有非零截面的点的轨迹在 $S$ 中闭。（把它施于外幂的扭，以及行列式的逆，以得到端点。）但这个轨迹恰好是 $|(BC(\mathcal{E})\setminus\{0\})/\underline{E}^\times|\to S$ 的像。由于这个映射紧合，像闭。

(2) 我们希望：若 $\mathrm{NP}(\mathcal{E})$ 为常值，则存在整体的 HN 滤过，且 pro-étale 局部有
\[
\mathcal{E}\cong\bigoplus_{\lambda\in\mathbb{Q}}\cO_{X_S}(\lambda)^{n_\lambda}.
\]
为此我们再次使用 v-下降。

我们断言，只需 v-局部地看到这一点。事实上，一旦有了它，就可以 v-局部地定义 HN 滤过，然后由唯一性它会下降（由 Fargues--Fontaine 曲线上向量丛的 v-下降，下次解释）。同构 $\mathcal{E}^\lambda\cong\cO_{X_S}(\lambda)^{n_\lambda}$ 构成 $\mathrm{GL}_{m_\lambda}(D_\lambda)$ 下的 v-主丛，由 [S17] 它必然来自 pro-étale 主丛。

\srcpage{48}

于是经过 pro-étale 局部化可以找到 $\mathcal{E}^\lambda\cong\cO_{X_S}(\lambda)^{n_\lambda}$。我们利用 $S$ 为 affinoid 时 $\lambda>0$ 的 $H^1(X_S,\cO_{X_S}(\lambda))=0$ 来分裂 HN 滤过。

剩下的任务是在 v-局部产生上述分裂。设 $\lambda$ 是 $\mathcal{E}$ 的最大斜率。我们想在一个 v-覆盖之后找到纤维上非零的映射 $\cO_{X_S}(\lambda)\to\mathcal{E}$。有了它，就会得到短正合列
\[
0\to\cO_{X_S}(\lambda)\to\mathcal{E}\to\overline{\mathcal{E}}\to0,
\]
再由归纳获胜。（仍然用到它 pro-étale 局部分裂。）

但现在令 $\mathcal{E}'=\underline{\Hom}(\cO_{X_S}(\lambda),\mathcal{E})$，我们有
\[
BC(\mathcal{E}')\setminus\{0\}\twoheadrightarrow
(BC(\mathcal{E}')\setminus\{0\})/\underline{E}^\times\to S,
\]
而后一个映射紧合，且在几何点上满射，因为我们在几何点上已经知道结论。$(BC(\mathcal{E}')\setminus\{0\})/\underline{E}^\times\to S$ 的紧合性（从而拟紧性）结合命题 9.8 说明它是 v-覆盖。映射 $BC(\mathcal{E}')\setminus\{0\}\twoheadrightarrow(BC(\mathcal{E}')\setminus\{0\})/\underline{E}^\times$ 按定义是 v-覆盖（它是 v-层的满射）。所以复合是 v-覆盖。但在 $BC(\mathcal{E}')\setminus\{0\}$ 上，我们重言地有一个纤维上非零的映射 $\cO_{X_S}(\lambda)\to\mathcal{E}$。\hfill$\square$

\begin{example}
$(BC(\cO(1))\setminus\{0\})/\underline{E}^\times$ 长什么样？几何上，$BC(\cO(1))\setminus\{0\}$ 是一个完满的去心开单位圆盘。$\pi$ 的作用向原点收缩，把半径沿其边界粘起来。

事实上，由 Lubin--Tate 理论有 $BC(\cO(1))\setminus\{0\}=(\Spa\breve E^{\mathrm{LT}}_\infty)^\diamond$，所以
\[
(BC(\cO(1))\setminus\{0\})/\underline{E}^\times
\cong(\Spa\breve E^{\mathrm{LT}}_\infty)^\diamond/\underline{E}^\times
\cong(\Spa\breve E)^\diamond/\varphi^{\mathbb{Z}}
=\Div^1_X.
\]
另一种看法是：$BC(\cO(1))\setminus\{0\}$ 参数化非零映射 $\cO\hookrightarrow\cO(1)$。其余核是 $\cO_D$，其中 $D$ 是某个 1 次除子，这正是 $\Div^1_X$ 所参数化的。除子 $D$ 在把映射乘以可逆函数（即 $\underline{E}^\times$）下不变。更一般地，
\[
(BC(\cO(d))\setminus\{0\})/\underline{E}^\times\cong\Div^d_X.
\]
\end{example}
